\documentclass[a4paper,12pt]{article}

\usepackage[utf8]{inputenc}
\usepackage[T1]{fontenc}
\usepackage[ngerman]{babel}
\usepackage{graphicx}
\usepackage{csquotes} 
\usepackage{enumitem}

% Da biblatex verwendet wird, bietet es sich an im Übersetzungsprozess
% biber statt bibtex zu verwenden
\usepackage[style=apa]{biblatex}

\usepackage{hyperref}
\usepackage{acronym} 
\renewcommand{\thesection}{\arabic{section}}
\newcounter{savepage}

\usepackage[left=25mm,right=30mm,top=25mm,bottom=25mm]{geometry}
\usepackage[onehalfspacing]{setspace}

\addto\captionsngerman{
  \renewcommand{\contentsname}
    {Gliederung}
}

\addbibresource{Quellen.bib}

\begin{document}
\begin{titlepage}
{{ \noindent  Universität Leipzig\\
		Wirtschaftswissenschaftliche Fakultät\\
		\itshape Betreuender Hochschullehrer\\
		Betreuender Assistent\\	 \par}
	\vspace{2cm}
	{\centering  Thema \\
	\bfseries \large [Thema der Arbeit]\par}
	\vspace{2cm}
	{ \centering  Masterarbeit zur Erlangung des akademischen Grades\\
	     Master of Science – [Studiengang] 
		\par}

	\vfill}

	{\noindent Vorgelegt von: \textit{Name, Vorname } \\
	 Matrikelnummer: \textit{XXXXXX} \\
	 Email-Adresse: \textit{adresse@mail.de} \\
	 Telefonnummer: \textit{XXXXXXXXXXXX} \\
	 Anschrift:  {Straße Hausnummer\\
						\hspace*{2.15cm} Postleitzahl Ort } \\ } 
		
		{\noindent  Leipzig, den Abgabedatum\par}
\end{titlepage}

\pagenumbering{Roman}

%Inhaltsverzeichnis
\addcontentsline{toc}{section}{Gliederung}
\tableofcontents
\newpage

%Abbildungsverzeichnis, bei Bedarf
\addcontentsline{toc}{section}{Abbildungsverzeichnis}
\listoffigures
\newpage

%Abkürzungsverzeichnis, bei Beadarf
\addcontentsline{toc}{section}{Abkürzungsverzeichnis}
{\noindent\Large\bfseries Abkürzungsverzeichnis}
\begin{acronym}
	\acro{GDP}{Gross domestic product} 
\end{acronym} 
\newpage

%Symbolverzeichnis, bei Bedarf
%\addcontentsline{toc}{section}{Symbolverzeichnis}
%{\noindent\Large\bfseries Symbolverzeichnis}
%\newpage

%Tabellenverzeichnis, bei Bedarf
\addcontentsline{toc}{section}{Tabellenverzeichnis}
\listoftables
\newpage
\clearpage

%Aktuelle römische Seitenzahl speichern. um später weiterzuführen
\setcounter{savepage}{\arabic{page}}

\pagenumbering{arabic}

%Include each chapter
\chapter{Einleitung}
Anstelle \textit{Ihrer} Einleitung steht hier zunächst eine kurze Einleitung in die Benutzung dieser Formatvorlage.

\subsection{PDF-Erstellung und eigener Inhalt}
Zur Übersetzung dieser Vorlage müssen Sie biber verwenden, da Biblatex für die Quellenverwaltung verwendet wurde. \\
Für Ihre eigenen Kapitel können Sie eine neue .tex-Datei im Unterordner Kapitel anlegen und diese in der Hauptdatei referenzieren. In Texmaker bietet es sich an die Vorlage\_ Latex als Masterdatei zu definieren. So kann auch in den Unterdateien die Funktion Schnelles Übersetzen verwendet werden.
Eigene Quellenverzeichnisse können Sie über einen Verweis in der Präambel einbinden.


\subsection{Zitationsbeispiele}
Zitate können Sie über die Bordmittel simpel einbinden, da diese automatisch an den gewählten Zitationsstil angepasst werden \parencite[5-8]{bergener_wissenschaftliches_2019}.
Auch Zitate im Fließtext \textcite{bergener_wissenschaftliches_2019} lassen sich abbilden. Die Grundlagen zur neuesten APA-Zitationsweise können Sie im Originalwerk nachlesen \parencite{apa_2019}, alternativ zur Harvard-Zitatsweise. Grundlage dieser Arbeit bildet die Harvard-Zitierweise.

\subsection{Tabellen und Bilder}

\begin{figure}[htbp] 
  \centering
     \includegraphics{Bilder/buchmitkopf.png}
  \caption[Buch mit Kopf]{Buch mit Kopf \parencite[1-2]{bergener_wissenschaftliches_2019}}
  \label{fig:Bild1}
\end{figure}

\begin{table}[htbp]
%Table-Umgebung erstellen
\centering
%Bei Harvard-Zitierweise kommen Beschriftungen über die Tabelle
\caption[Beschreibung der Tabelle]{Beschreibung der Tabelle \parencite[1-2]{bergener_wissenschaftliches_2019}}
%Tabelle erstellen
\begin{tabular}{|p{0.3\textwidth}|p{.3\textwidth}|p{.3\textwidth}|}
\hline
\textbf{Überschrift} & \textbf{Überschrift} & \textbf{Überschrift} \\ \hline
Zelltext             & Zelltext             &   Zelltext             \\ \hline
Zelltext             & Zelltext             & Zelltext             \\ \hline
\end{tabular}
\label{tab:beispieltabelle}
\end{table}
Im Gegensatz zu Word kann in Latex der Text für Tabellenverzeichnis und Tabellenunterschrift unterschiedlich definiert werden. Es bietet sich daher an die Zitation zur Tabelle direkt in der Unterschrift zu tätigen, statt dies im Fließtext nachzuholen. Ein Beispiel für die Verwendung von Tabellen sehen Sie mit Tabelle \ref{tab:beispieltabelle}. Verschiedene Tabellengeneratoren im Internet erleichtern die Verwendung von Tabellen enorm.












% Literaturverzeichnis
\newpage
\pagenumbering{Roman}
\setcounter{page}{\thesavepage}
\addcontentsline{toc}{section}{Literaturverzeichnis}
{\noindent\Large\bfseries Literaturverzeichnis}
\singlespacing

\printbibliography
\onehalfspacing
%Anhang, bei Bedarf
\newpage
\addcontentsline{toc}{section}{Anhang}
{\noindent\Large\bfseries Anhang}


\end{document}